\documentclass[11pt]{article}
\usepackage[margin=1in]{geometry}
\usepackage{amsmath}
\usepackage{booktabs}
\usepackage{hyperref}

\title{NLP Assignment 3}
\author{Adam Fleisher \and Gal Barak \and Tamar Tabbach}
\date{}

\begin{document}
\maketitle

\section*{Question 1: Prompting for Reasoning}

% ============================================================
\subsection*{1.2.d -- Zero-Shot Accuracies and Model Comparison}
% ============================================================

\begin{table}[h]
\centering
\begin{tabular}{lc}
\toprule
\textbf{Method} & \textbf{Accuracy} \\
\midrule
LLaMA zero-shot (temperature = 0) & 0.68 \\
GPT-2 zero-shot (temperature = 0) & 0.48 \\
\bottomrule
\end{tabular}
\end{table}

LLaMA outperforms GPT-2 by 20 percentage points (0.68 vs.\ 0.48). Three differences that may explain this gap:

\begin{enumerate}
    \item \textbf{Model scale and training.} LLaMA-3.1-8B is a modern instruction-tuned model trained on large-scale data for general reasoning and QA. GPT-2 (124M parameters for the hugging face version we are loading) is a much smaller, older autoregressive language model trained mainly on next-token prediction without instruction tuning.
    \item \textbf{Prompting interface.} LLaMA is queried through a chat API with an explicit system prompt instructing the model to answer only \texttt{True} or \texttt{False}. GPT-2 has no system role; the instruction must be embedded in the raw text prompt, which is a weaker way to enforce output format, as we can see in the notebook a lot of these answers are neither correct or incorrect they are completely not in the format and therefore invalid (but marked incorrect).
    \item \textbf{Task suitability.} StrategyQA requires multi-hop commonsense reasoning. LLaMA was trained on diverse instruction-following and reasoning data, while GPT-2 lacks this specialization and often generates free-form continuations rather than concise boolean answers.
\end{enumerate}

\newpage
\subsection*{1.3.b -- Self-Consistency Accuracy (temperature = 1.0)}

Using self-consistency with $n = 5$ samples and temperature $= 1.0$, the accuracy was: \textbf{0.64}

\subsection*{1.3.c -- Temperature Analysis (0, 0.5, 1.0)}

We repeated self-consistency ($n=5$) at three temperature settings:

\begin{center}
\begin{tabular}{lccc}
\toprule
\textbf{Temperature} & \textbf{Accuracy} & \textbf{Unanimous votes} & \textbf{Split votes} \\
\midrule
0.0  & 0.680 & 50/50 (100\%) & 0/50 (0\%) \\
0.5  & 0.640 & 40/50 (80\%)  & 10/50 (20\%) \\
1.0  & 0.640 & 37/50 (74\%)  & 13/50 (26\%) \\
\bottomrule
\end{tabular}
\end{center}

\textbf{Vote distributions (selected patterns):}
\begin{itemize}
    \item \textbf{Temp = 0:} Only two patterns appeared: \texttt{True:0, False:5} (37 questions) and \texttt{True:5, False:0} (13 questions). Every question had identical votes across all 5 samples. 
    This is because when temperature is set to 0, the model just takes argmax over its answer distribution making it deterministic, this means that self-consistency is meaningless when temperature is set to 0.
    
    \item \textbf{Temp = 0.5:} Here sampling adds randomness, but the distribution stays sharper than at temp 1.0, so high-probability answers are still chosen most of the time — hence mostly 5–0 votes, with occasional splits on uncertain questions.
    
    \item \textbf{Temp = 1.0:} Highest diversity: 13/50 questions had split votes.
\end{itemize}

\textbf{Insights:}
\begin{itemize}
    \item \textbf{Temperature vs.\ diversity:} Lower temperature produces nearly deterministic outputs (100\% unanimous at $T=0$). Higher temperature increases output diversity and the frequency of split votes (26\% at $T=1.0$).
    \item \textbf{Self-consistency benefit:} In our case, since when $T=0$ is equal to no majority voting at all, it seems that self consistency only damaged the results. This of-course does not prove anything since the experiment scope is relatively small and running the same experiments with a larger dataset or with different random seeds might provide different results.
\end{itemize}


\newpage
\subsection*{1.4.d -- ICL Accuracy ($n = 10$)}

Using in-context learning with $n = 10$ demonstrations (10 random examples from \texttt{demos\_dict}), temperature $= 0$: \textbf{0.62}

\subsubsection*{1.4.e --- ICL Analysis}

We analyzed how demonstration \textbf{quantity}, \textbf{ordering}, and \textbf{similarity} affect accuracy.
All experiments used temperature $0$ and the same evaluation pipeline as in 1.4.d. \textbf{Important Note:} Since the examples are chosen randomly, rerunning the experiment showed high variance in the results, we chose a representing experiment to base our conclusions on (avaraging over many experiments takes considerable time with our computers).

\paragraph{Quantity of demonstrations.}
\begin{center}
\begin{tabular}{lc}
\hline
Number of demos ($n$) & Accuracy \\
\hline
1  & 0.580 \\
3  & 0.600 \\
5  & 0.580 \\
10 & 0.600 \\
20 & 0.580 \\
\hline
\end{tabular}
\end{center}

Accuracy was relatively unstable and did not increase monotonically with more demonstrations.
The best result in this run was $n=3$ and $n=10$ (both 0.60), while $n=1$, $n=5$, and $n=20$ were slightly lower (0.58).
This suggests that more examples do not automatically help: very long prompts may dilute attention, and additional demos may add noise rather than useful signal. In our runs, performance stayed in a narrow band (0.58--0.60) for most values of $n$, with no clear benefit from scaling up the number of demonstrations.

\paragraph{Ordering of demonstrations ($n=10$).}
To isolate ordering effects, we first sampled one fixed set of 10 demonstrations, then tested three orderings of the \emph{same} demos:
\begin{center}
\begin{tabular}{lc}
\hline
Ordering & Accuracy \\
\hline
Original order  & 0.640 \\
Reversed order  & 0.680 \\
Shuffled order  & 0.540 \\
\hline
\end{tabular}
\end{center}

Reversing the demonstrations improved accuracy slightly (0.68), while shuffling hurt (0.54).
This indicates that order can matter, likely because recent/context-adjacent examples may influence the model more strongly.
However, the effect is not fully consistent, so ordering alone does not explain ICL performance.

\paragraph{Similarity of demonstrations ($n=10$).}
We compared:
(i) 10 random demos per question, and
(ii) 10 demos selected per question by highest word-overlap similarity (Jaccard similarity over question tokens).
\begin{center}
\begin{tabular}{lc}
\hline
Demo selection strategy & Accuracy \\
\hline
Random 10 demos                   & 0.600 \\
10 most similar demos per question & 0.540 \\
\hline
\end{tabular}
\end{center}

Similarity-based retrieval did \emph{not} improve performance.
Simple lexical overlap is a weak proxy for reasoning relevance on StrategyQA, so ``similar-looking'' questions may still teach the wrong reasoning pattern. This result initially surprised us but then we thought about it being some sort of overfitting, questions with overlapping word does not necessarily mean the same answer, especially in yes or no questions...

\paragraph{Summary.}
In our experiments, ICL gains were limited and sensitive to demo configuration.
The strongest setting was reversed order with a fixed strong random subset (0.68), matching zero-shot, while random shuffling and similarity-based selection performed worse.
Overall, ICL only helps when demonstrations are informative and well presented, otherwise it can underperform zero-shot.

\newpage
\subsection*{1.5.c -- Chain-of-Thought Accuracy}
Using chain-of-thought prompting with a custom system prompt, temperature $= 0$, and \texttt{max\_tokens = 1024}: \textbf{0.76}

CoT improved accuracy from 0.68 (zero-shot) to 0.76 (+8 points), showing that explicit step-by-step reasoning helps on multi-hop StrategyQA questions.

% ============================================================
\subsection*{1.6.d -- ICL + CoT Combined Accuracy}
% ============================================================

Using 1, 3, and 5 chain-of-thought demonstrations (with facts and decomposition) plus CoT prompting on the target question we got: \textbf{0.54, 0.60, 0.60}
meaning 3 and 5 chain-of-thought demonstrations tied for the best results in this run.

\subsubsection*{1.6.e --- Method Comparison}

We compared all prompting methods on the 50-question StrategyQA test set using Llama~3.1-8B with temperature $0$:

\begin{center}
\begin{tabular}{lc}
\hline
Method & Accuracy \\
\hline
Zero-shot          & 0.68 \\
ICL ($n=10$)       & 0.62 \\
CoT                & \textbf{0.76} \\
ICL + CoT ($n=1$)  & 0.54 \\
ICL + CoT ($n=3$)  & 0.60 \\
ICL + CoT ($n=5$)  & 0.60 \\
\hline
\end{tabular}
\end{center}

\paragraph{Discussion.}
Plain \textbf{ICL} was the weakest method among the main baselines (0.62), below zero-shot (0.68): providing only question--answer pairs, without any reasoning, gives the model a pattern to copy but no guidance on \emph{how} to solve multi-hop questions.

Adding reasoning helps: \textbf{CoT} reached \textbf{0.76}, the best overall result, since letting the model think step by step before answering fits the multi-step nature of StrategyQA.

For \textbf{ICL + CoT} we varied the number of demonstrations with a fixed random demo set per setting. Performance was again sensitive to this choice, but it did \emph{not} beat pure CoT: $n=1$ reached only 0.54, while $n=3$ and $n=5$ both reached 0.60. Combining demonstrations with CoT therefore did not improve over CoT alone in this final run; the longer prompts likely added noise and made answer extraction harder, while a single demonstration was especially unhelpful.

\paragraph{Conclusion.}
Reasoning-based prompting clearly outperforms plain ICL, and \textbf{pure CoT (0.76)} was the strongest method overall. ICL + CoT remained unstable and underperformed CoT in our final results, suggesting that adding solved examples is not automatically beneficial once the model is already prompted to reason step by step. The method also appears sensitive to the chosen demonstrations, prompt length, and API runtime conditions, so a more comprehensive experiment averaged over multiple demo sets would be needed to draw stronger conclusions.

\end{document}